\documentclass{beamer}
\usepackage[utf8]{inputenc}
\usetheme{Madrid}
\usecolortheme{default}
\usepackage{amsmath,amssymb,amsfonts,amsthm}
\usepackage{txfonts}
\usepackage{tkz-euclide}
\usepackage{listings}
\usepackage{adjustbox}
\usepackage{array}
\usepackage{tabularx}
\usepackage{gvv}
\usepackage{lmodern}
\usepackage{circuitikz}
\usepackage{tikz}
\usepackage{graphicx}

\setbeamertemplate{page number in head/foot}[totalframenumber]

\usepackage{tcolorbox}
\tcbuselibrary{minted,breakable,xparse,skins}



\definecolor{bg}{gray}{0.95}
\DeclareTCBListing{mintedbox}{O{}m!O{}}{%
  breakable=true,
  listing engine=minted,
  listing only,
  minted language=#2,
  minted style=default,
  minted options={%
    linenos,
    gobble=0,
    breaklines=true,
    breakafter=,,
    fontsize=\small,
    numbersep=8pt,
    #1},
  boxsep=0pt,
  left skip=0pt,
  right skip=0pt,
  left=25pt,
  right=0pt,
  top=3pt,
  bottom=3pt,
  arc=5pt,
  leftrule=0pt,
  rightrule=0pt,
  bottomrule=2pt,
  toprule=2pt,
  colback=bg,
  colframe=orange!70,
  enhanced,
  overlay={%
    \begin{tcbclipinterior}
    \fill[orange!20!white] (frame.south west) rectangle ([xshift=20pt]frame.north west);
    \end{tcbclipinterior}},
  #3,
}
\lstset{
    language=C,
    basicstyle=\ttfamily\small,
    keywordstyle=\color{blue},
    stringstyle=\color{orange},
    commentstyle=\color{green!60!black},
    numbers=left,
    numberstyle=\tiny\color{gray},
    breaklines=true,
    showstringspaces=false,
}
%------------------------------------------------------------
%This block of code defines the information to appear in the
%Title page
\title %optional
{5.13.2}

%\subtitle{A short story}

\author % (optional)
{BALU-ai25btech11017}



\begin{document}


\frame{\titlepage}
\begin{frame}{Question}
If $\vec{A}=\myvec{a&&b\\b&&a}$ and $\vec{A}^2=\myvec{\alpha&&\beta\\\beta&&\alpha}$,then
\\ 
\end{frame}
\begin{frame}{Theoretical Solution}
Let us solve the given equation theoretically and then verify the solution computationally \\
According to the question, \\
We known that\\
\begin{align}
    \|\vec{A}-\lambda\vec{I}\|=0\\
    \left\|\myvec{a-\lambda&&b\\b&&a-\lambda}\right\|=0\\
    \lambda^2-2a\lambda+a^2-b^2=0\
\end{align}
using Cayley-Hamilton theorem
\end{frame}
\begin{frame}{Theoretical Solution}
\begin{align}
    \vec{A}^2-2a\vec{A}+(a^2-b^2)\vec{I}=0\\
    \vec{A}^2=2a\vec{A}-(a^2-b^2)\vec{I}\\
    \vec{A}^2=\myvec{a^2+b^2&&2ab\\2ab&&a^2+b^2}\\
    \alpha=a^2+b^2,\beta=2ab
\end{align}
\end{frame}
\begin{frame}[fragile]
    \frametitle{C Code}

    \begin{lstlisting}
  #include<stdio.h>
int main(){
    int arr[2][2];
    for(int i=0;i<2;i++){
        for(int j=0;j<2;j++){
            scanf("%d",&arr[i][j]);
        }
    }
    int arr2[2][2]={0};
    for(int i=0;i<2;i++){
        for(int j=0;j<2;j++){
            for(int z=0;z<2;z++){
                arr2[i][j]+=arr[i][z]*arr[z][j];
            }
        }
    }
    
    \end{lstlisting}
\end{frame}
\begin{frame}[fragile]
    \frametitle{C Code }
    \begin{lstlisting}
for(int i=0;i<2;i++){
        for(int j=0;j<2;j++){
            printf("%d",arr2[i][j]);
        }
        printf("\n");
    }
    return 0;
}
   
    \end{lstlisting}
\end{frame}






\end{document}